% !TEX root = ../../cookbook.tex
\newpage
\recipe{Engadine Walnut Tart}{}{}{}
\label{bundner_nusstorte}

\begin{multicols}{2}
	\subsection*{Ingredients}
	\subsubsection*{Dough}
	\begin{ingredients}
		\item 325 g flour
		\item 100 g sugar
		\item 1 tsp baking powder
		\item 1 pinch salt
		\item 1 egg
		\item 2 tbsp cold water
		\item 175 g butter
	\end{ingredients}
	
	\columnbreak
	
	\subsection*{}
	\subsubsection*{Filling}
	\begin{equipment}
		\item 200 g sugar
		\item 0.5 dl hot water
		\item 2.5 dl heavy cream
		\item 300 g walnuts
		\item 50 g honey
	\end{equipment}
\end{multicols}

\medskip

% --- Preparation ---
Mix the \textbf{flour} and \textbf{butter} using a mixer.
Then combine the flour-butter mixture with the \textbf{sugar}, \textbf{baking powder}, and \textbf{salt} in a bowl.
Add the \textbf{egg} and the \textbf{water}, and mix with a fork.
Wrap in plastic wrap and refrigerate for about 30 minutes. \\

For the filling, caramelize the \textbf{sugar} \underline{slowly}, pour in the \textbf{water}, and bring to a boil.
Add the \textbf{cream} and simmer until the caramel has dissolved.
Stir in the \textbf{walnuts} and \textbf{honey}.
Continue stirring for a few minutes until the mixture thickens.
Then let it cool. \\

Preheat the oven to 200°C conventional heat (top and bottom).
Divide the dough into three parts.
Roll out one third into a disc about 24 cm in \diameter\ and refrigerate.
Roll out the remaining dough on a lightly floured surface into a disc about 32 cm in \diameter\ and place it in the buttered pan.
Spread the filling evenly over the dough.
Fold the edges of the dough inward and cover the tart with the smaller dough disc.
Press the edges firmly with a fork and prick the top.
Bake on the middle rack for about 30 minutes.
Allow to cool completely (this is quite important).

\vspace*{0.75 cm}

% --- Description ---
\begin{recipeblock}
	\noindent The \textit{Engadiner} (or \textit{Bündner}) \textit{Nusstorte} (`nut pie') is traditional of the alpine region of \textit{Engiadina}, in the Kanton of Graubünden, Switzerland. This is my favorite pie ever. 
\end{recipeblock}






